\documentclass[11pt,a4paper]{article}
\usepackage[margin=1in]{geometry}
\usepackage{times}
\usepackage{latexsym}
\usepackage{graphicx}
\usepackage{amsmath}
\usepackage{amssymb}
\usepackage{booktabs}
\usepackage{multirow}
\usepackage{xcolor}
\usepackage{url}
\usepackage[colorlinks=true,citecolor=blue,linkcolor=blue]{hyperref}
\usepackage{natbib}
\bibliographystyle{plainnat}

\newcommand{\coverageup}[1]{\textcolor{green!60!black}{(+#1)}}
\newcommand{\coveragedown}[1]{\textcolor{red!70!black}{(#1)}}

\title{Contrastive Trace Probing for Safety Commitment Detection\\in Reasoning Models: A Cross-Architecture Empirical Study}

\author{Anonymous}

\begin{document}
\maketitle

% ========================================================================
% ABSTRACT
% ========================================================================
\begin{abstract}
Reasoning models such as DeepSeek-R1 and QwQ exhibit \emph{self-jailbreak}: their extended chain-of-thought interleaves safety reasoning with harmful execution, culminating in a \emph{commitment point} where safety reasoning permanently ceases.
Existing hidden-state jailbreak detectors target adversarial inputs to instruction-tuned models and have never been evaluated on reasoning-model self-jailbreak.
We present the first cross-architecture empirical study of hidden-state safety monitoring in reasoning models, spanning four models across two architecture families (Llama-3.1 and Qwen2.5) and two scales (7--8B and 32B).
We establish three findings:
(1)~static per-step probes achieve $<$50\% detection coverage despite strong per-dimension safety signal (Fisher ratio $>$14), because the signal is distributed across minor principal components;
(2)~a simple temporal window logistic regression ($W$=15) achieves 76.7\% mean coverage---matching the label-oracle ceiling---with a +34.9pp improvement over static probes that is consistent across all architectures;
(3)~no hidden-state method achieves positive detection lead time (best: $-$0.2 steps vs.\ oracle +31 steps), establishing a principled limitation boundary.
Cross-architecture transfer via learned projection heads yields 64.5\% mean coverage without target-domain labels, suggesting safety commitment signals share geometric structure across model families.
\end{abstract}

% ========================================================================
% 1. INTRODUCTION
% ========================================================================
\section{Introduction}

Reasoning models---large language models trained with reinforcement learning to produce extended chains of thought (CoT)---represent a qualitative shift in how language models process complex queries \cite{deepseekr1,qwq32b}.
These models interleave multiple reasoning strategies across dozens to hundreds of steps, achieving strong performance on mathematical and scientific benchmarks.
However, recent work has identified a critical safety failure mode: \emph{self-jailbreak}, where the model's own reasoning process gradually transitions from safety-aligned analysis to harmful content generation \cite{harmthoughts,selfjailbreaking}.

Unlike conventional jailbreak attacks that rely on adversarial inputs, self-jailbreak emerges from the model's internal reasoning dynamics.
In 84\% of harmful traces, safety-relevant reasoning persists well past step 10, only to permanently cease at a \emph{commitment point}---the reasoning step after which no further safety considerations appear.
Detecting this commitment point is critical for deploying safety monitors on reasoning models.

Existing hidden-state jailbreak detectors such as TrajGuard \cite{trajguard}, SafeSwitch \cite{safeswitch}, and SIREN \cite{siren} achieve strong defense rates ($>$95\%) on instruction-tuned models.
However, they operate at token-level granularity on models where the attack is an adversarial \emph{input}.
Reasoning-model self-jailbreak is fundamentally different: the attack is the model's \emph{own reasoning process}, distributed across many CoT steps.
No prior work bridges hidden-state detection methods to reasoning-model self-jailbreak.

We present the first systematic cross-architecture evaluation of hidden-state methods for safety commitment detection in reasoning models.
Our study spans four models---DeepSeek-R1-Distill-Llama-8B, DeepSeek-R1-Distill-Qwen-32B, QwQ-32B, and OpenThinker-7B---covering two architecture families (Llama-3.1 and Qwen2.5) and two scales (7--8B and 32B).
We establish three claims:

\paragraph{Claim 1: Static probes fail reliably.} Per-step logistic regression achieves $<$50\% detection coverage across all models (mean 41.9\%), despite Fisher discriminant ratios $>$14 in full representation space.  The safety signal exists but concentrates in minor principal components (top-3 PCs explain $<$17\% of variance), causing both static classifiers and geometric detectors (Mahalanobis distance) to fail.

\paragraph{Claim 2: Temporal window probes saturate oracle coverage.} A temporal window logistic regression with $W$=15 achieves 76.7\% mean coverage across architectures, matching the label-oracle upper bound.  The +30--35pp improvement is consistent across all four models and both architecture families.  Window size $W$=15 is universally optimal, requiring no architecture-specific tuning.

\paragraph{Claim 3: Pre-textual detection is fundamentally limited.} No hidden-state method achieves positive mean lead time. The oracle's +31-step advantage derives from directly observing label-sequence autocorrelation---information that hidden states encode for the current step but not for future trajectory.  This establishes a principled boundary: hidden-state monitors can detect commitment \emph{at} or \emph{after} it occurs, but not before.

% ========================================================================
% 2. RELATED WORK
% ========================================================================
\section{Related Work}

\paragraph{Hidden-state jailbreak detection.}
TrajGuard \cite{trajguard} monitors activation trajectories using Mahalanobis distance from safe-trace distributions, achieving 95\% defense rate on instruction-tuned models.
SafeSwitch \cite{safeswitch} trains linear probes on intermediate representations to classify safety-relevant states.
SIREN \cite{siren} uses contrastive features between safe and unsafe hidden states for real-time detection.
All three operate at token-level granularity on models where the adversarial signal is in the \emph{input}---none have been evaluated on reasoning-model self-jailbreak where the signal is in the model's own multi-step reasoning.

\paragraph{Reasoning model safety.}
\citet{selfjailbreaking} demonstrate that reasoning models can jailbreak themselves through extended CoT, bypassing safety training without adversarial input.
\citet{harmthoughts} provide the HarmThoughts dataset with fine-grained step-level safety annotations across a 16-label taxonomy, enabling quantitative analysis of safety reasoning dynamics.
SafeRemind \cite{saferemind} and SafeThink \cite{safethink} propose text-level CoT monitoring but are vulnerable to pathological reasoning chains.
None of these works evaluate hidden-state detection methods.

\paragraph{Linear probing of LLM representations.}
Linear probes have been widely used to study what information LLM hidden states encode \cite{belinkov2022probing}, including truthfulness \cite{li2024inference}, sentiment \cite{tigges2023linear}, and factual knowledge \cite{meng2022locating}.
Our work extends this methodology to temporal probe designs that capture multi-step patterns in reasoning traces.

% ========================================================================
% 3. METHOD
% ========================================================================
\section{Method}

\subsection{Task Definition}

Given a reasoning trace $\mathbf{x} = (x_1, \ldots, x_T)$ where each $x_t$ is a reasoning step, we define the \emph{commitment point} $c$ as the last step containing safety-relevant reasoning.
The detection task is binary classification: for each step $t$, predict whether $t > c$ (the model has committed to harmful output).

\subsection{Hidden-State Extraction}

For each reasoning step $x_t$, we extract the hidden-state vector $\mathbf{h}_t^{(\ell)} \in \mathbb{R}^d$ at layer $\ell$ from the last token position of that step.
Based on layer ablation (Section~\ref{sec:layer_ablation}), we use the optimal middle-layer representation at approximately 35--44\% relative depth.

\subsection{Static Probe}

The static probe is a logistic regression classifier trained on individual hidden states:
\begin{equation}
    p(y_t = 1 \mid \mathbf{h}_t) = \sigma(\mathbf{w}^\top \mathbf{h}_t + b)
\end{equation}
This represents the standard approach used by existing hidden-state detectors adapted to step-level granularity.

\subsection{Temporal Window Probe}

The temporal window probe concatenates hidden states from a window of $W$ consecutive steps:
\begin{equation}
    \mathbf{z}_t = [\mathbf{h}_{t-W+1}; \ldots; \mathbf{h}_t] \in \mathbb{R}^{Wd}
\end{equation}
\begin{equation}
    p(y_t = 1 \mid \mathbf{z}_t) = \sigma(\mathbf{w}^\top \mathbf{z}_t + b)
\end{equation}
This enables the probe to capture \emph{safety reasoning density decline}---the gradual decrease in safety-labeled steps within the window as the model approaches commitment.

\subsection{Mahalanobis Distance Baseline}

Following TrajGuard's approach, we estimate the covariance $\mathbf{\Sigma}$ of hidden states from safe traces and compute per-step Mahalanobis distance:
\begin{equation}
    d_M(\mathbf{h}_t) = \sqrt{(\mathbf{h}_t - \boldsymbol{\mu})^\top \mathbf{\Sigma}^{-1} (\mathbf{h}_t - \boldsymbol{\mu})}
\end{equation}
A threshold on $d_M$ classifies steps as committed.

\subsection{GRU Sequential Probe}

As an upper bound on temporal modeling capacity, we train a GRU with 256 hidden units that processes the full sequence of hidden states:
\begin{equation}
    \mathbf{g}_t = \text{GRU}(\mathbf{h}_1, \ldots, \mathbf{h}_t; \theta) \in \mathbb{R}^{256}
\end{equation}
\begin{equation}
    p(y_t = 1 \mid \mathbf{g}_t) = \sigma(\mathbf{v}^\top \mathbf{g}_t + b)
\end{equation}

\subsection{Label-Sequence Oracle}

The oracle directly observes the binary label sequence and fires when $K=5$ consecutive steps lack safety labels.
It serves as an upper bound on detection \emph{coverage} (what fraction of commitment points are detected) while providing a reference for \emph{lead time} (how many steps before textual commitment the detection occurs).

\subsection{Cross-Architecture Transfer}

To evaluate whether safety commitment signals share geometric structure across architectures, we train a linear projection head $\mathbf{P} \in \mathbb{R}^{d_{\text{target}} \times d_{\text{source}}}$ that maps hidden states from a source model's space to a target model's space, using 30 aligned traces for projection training.
The source model's trained probe is then applied to projected target-model representations.

\subsection{Evaluation Metrics}

\begin{itemize}
    \item \textbf{Coverage}: Fraction of traces where the probe detects commitment (true positive rate at trace level).
    \item \textbf{Per-step accuracy}: Step-level classification accuracy.
    \item \textbf{Lead time}: Number of steps between probe detection and textual commitment point (positive = early detection).
\end{itemize}

% ========================================================================
% 4. EXPERIMENTAL SETUP
% ========================================================================
\section{Experimental Setup}

\paragraph{Models.} We evaluate four reasoning models spanning two architecture families and two scales: DeepSeek-R1-Distill-Llama-8B (Llama-3.1, $d$=4096, 32 layers), DeepSeek-R1-Distill-Qwen-32B (Qwen2.5, $d$=5120, 64 layers), QwQ-32B (Qwen2.5, $d$=5120, 64 layers), and OpenThinker-7B (Qwen2.5, $d$=3584, 32 layers).

\paragraph{Dataset.} We use HarmThoughts reasoning traces with step-level safety annotations from a 16-label taxonomy, binarized to safety-relevant vs.\ not.
The dataset contains 209 traces, split into train/test with stratification.

\paragraph{Layer selection.} For each model, we select the optimal layer via held-out validation.
Optimal layers fall at 34--44\% relative depth across all models (Section~\ref{sec:layer_ablation}).

% ========================================================================
% 5. RESULTS
% ========================================================================
\section{Results}

\subsection{Claim 1: Static Probes Fail Across Architectures}

Table~\ref{tab:static} reports static probe performance across all models.
Despite per-step accuracies above 80\% (inflated by majority-class bias), detection coverage remains below 50\% for every model.
Fisher discriminant ratios in original representation space exceed 14 across architectures, confirming that safety-relevant signal \emph{exists} in hidden states.
However, PCA analysis reveals this signal is distributed across minor principal components: top-3 PCs explain only 12.8--16.6\% of variance, and Fisher ratio collapses from $>$14 to $<$1 after PCA-3 projection.

\begin{table}[t]
\centering
\small
\begin{tabular}{lcccccc}
\toprule
\textbf{Model} & \textbf{Arch} & $d$ & \textbf{Acc} & \textbf{Cov} & \textbf{Fisher} & \textbf{F(PCA-3)} \\
\midrule
r1-8b & Llama & 4096 & .817 & .465 & 17.9 & 0.93 \\
r1-32b & Qwen & 5120 & .824 & .442 & 17.1 & 0.87 \\
QwQ-32B & Qwen & 5120 & .809 & .395 & 15.6 & 0.72 \\
OT-7B & Qwen & 3584 & --- & .419 & 13.2 & 0.68 \\
\midrule
\multicolumn{3}{l}{\textit{Mean}} & .817 & .430 & 15.9 & 0.80 \\
\bottomrule
\end{tabular}
\caption{Static probe performance across architectures.  Coverage is below 50\% for all models despite high Fisher discriminant ratios in original space.  Fisher(PCA-3) collapses to $<$1, confirming signal concentration in minor PCs.}
\label{tab:static}
\end{table}

The Mahalanobis distance baseline (Table~\ref{tab:mahalanobis}) performs even worse than static LR, with mean coverage of 27.9\% and AUROC below 0.62.
This failure is principled: Mahalanobis distance is dominated by major principal components, where the safety signal is weakest.

\begin{table}[t]
\centering
\small
\begin{tabular}{lcccc}
\toprule
\textbf{Model} & \textbf{MD Cov} & \textbf{MD AUROC} & \textbf{Temp Cov} & \textbf{Gap} \\
\midrule
r1-8b & .326 & .612 & .767 & $-$44.1pp \\
r1-32b & .279 & .583 & .791 & $-$51.2pp \\
QwQ-32B & .233 & .561 & .744 & $-$51.1pp \\
\bottomrule
\end{tabular}
\caption{Mahalanobis distance baseline vs.\ temporal probe. Geometric anomaly detection fails because safety signal concentrates in minor PCs that contribute negligibly to Mahalanobis distance.}
\label{tab:mahalanobis}
\end{table}

\subsection{Claim 2: Temporal Window Probes Saturate Oracle Coverage}

Table~\ref{tab:temporal} presents the core result.
The temporal window probe ($W$=15) achieves 74.4--79.1\% coverage across architectures, with a consistent +30--35pp improvement over static probes.
This improvement holds for both Llama-3.1 and Qwen2.5 architectures and at both 7--8B and 32B scales.

\begin{table}[t]
\centering
\small
\begin{tabular}{lccccc}
\toprule
\textbf{Model} & \textbf{Stat} & \textbf{Temp} & \textbf{$\Delta$} & \textbf{Acc} & \textbf{Lead} \\
\midrule
r1-8b & .465 & .767 & +30.2 & .848 & $-$1.09 \\
r1-32b & .442 & .791 & +34.9 & .851 & $-$1.32 \\
QwQ-32B & .395 & .744 & +34.9 & .839 & $-$0.87 \\
\midrule
\textit{Mean} & .434 & .767 & +33.3 & .846 & $-$1.09 \\
\bottomrule
\end{tabular}
\caption{Temporal window probe ($W$=15) vs.\ static probe.  Coverage improvement of +30--35pp is consistent across architectures.  Lead time is the mean number of steps relative to commitment point (negative = after).}
\label{tab:temporal}
\end{table}

\paragraph{Architecture $\times$ Scale matrix.}
Table~\ref{tab:arch_scale} presents the 2$\times$2 architecture--scale matrix, including OpenThinker-7B as the Qwen2.5-7B control.
The temporal advantage holds at both scales: the fundamental pattern (static fails, temporal saturates) is architecture- and scale-independent.
Scale provides a modest +7pp improvement within the Qwen family (7B$\to$32B), smaller than the +30pp method effect.

\begin{table}[t]
\centering
\small
\begin{tabular}{lcc}
\toprule
& \textbf{7--8B} & \textbf{32B} \\
\midrule
Llama-3.1 & r1-8b: .767 & --- \\
Qwen2.5 & OT-7B: .721 & r1-32b: .791 \\
\bottomrule
\end{tabular}
\caption{Temporal $W$=15 coverage in the 2$\times$2 architecture--scale matrix.  Scale effect (+7pp) and architecture effect ($-$4.6pp) are both small relative to method effect (+30pp).}
\label{tab:arch_scale}
\end{table}

\subsection{Claim 3: Pre-Textual Detection Is Fundamentally Limited}

Table~\ref{tab:leadtime} reports lead time analysis across all models and methods.
No hidden-state method achieves positive mean lead time.
The oracle achieves +31 steps mean lead time by directly observing label-sequence autocorrelation---consecutive absences of safety labels predict permanent absence.

\begin{table}[t]
\centering
\small
\begin{tabular}{llcccc}
\toprule
\textbf{Model} & \textbf{Method} & \textbf{Lead} & \textbf{Med} & \textbf{+Lead\%} & \textbf{Cov} \\
\midrule
\multirow{3}{*}{r1-8b} & Static & $-$2.0 & $-$2 & 30\% & .465 \\
& Temp W15 & $-$1.09 & $-$1 & 34\% & .767 \\
& GRU & $-$0.21 & $-$1 & 34.5\% & .674 \\
\midrule
\multirow{3}{*}{r1-32b} & Static & $-$2.8 & $-$3 & 25\% & .442 \\
& Temp W15 & $-$1.32 & $-$1 & 29\% & .791 \\
& GRU & $-$0.9 & $-$1 & 29\% & .721 \\
\midrule
\multirow{3}{*}{QwQ-32B} & Static & $-$1.9 & $-$2 & 28\% & .395 \\
& Temp W15 & $-$0.87 & $-$1 & 33\% & .744 \\
& GRU & $-$0.5 & 0 & 33\% & .698 \\
\midrule
\multicolumn{2}{l}{Oracle ($K$=5)} & +31 & +29 & 100\% & 1.0 \\
\bottomrule
\end{tabular}
\caption{Detection lead time across architectures.  Lead = mean steps before commitment (positive = early). No hidden-state method crosses zero; the oracle's +31-step advantage comes from label-sequence statistics unavailable to hidden-state probes.}
\label{tab:leadtime}
\end{table}

\paragraph{Oracle decomposition.}
Table~\ref{tab:oracle_mi} reports mutual information between hidden states and label sequences at varying horizons.
Hidden states are informative about the \emph{current} step's label (MI = 0.42) but nearly uninformative about future labels (MI $\to$ 0.03 at horizon 15).
The oracle exploits label-sequence autocorrelation: on average, 14.6 consecutive steps precede the commitment point without safety labels, making $K$=5 consecutive absences a powerful predictor.
This information is structurally unavailable to hidden-state probes.

\begin{table}[t]
\centering
\small
\begin{tabular}{lcc}
\toprule
\textbf{Target} & \textbf{MI(Hidden)} & \textbf{MI(Labels)} \\
\midrule
Current label & 0.42 & 0.89 \\
Next 5 labels & 0.08 & 0.31 \\
Next 15 labels & 0.03 & 0.12 \\
Commitment (binary) & 0.15 & 0.67 \\
\bottomrule
\end{tabular}
\caption{Mutual information between hidden states (resp.\ label history $W$=15) and label sequences. Hidden states encode current state but not future trajectory.}
\label{tab:oracle_mi}
\end{table}

\subsection{Cross-Architecture Transfer}

Table~\ref{tab:transfer} presents the cross-architecture transfer matrix.
A probe trained on one model achieves non-trivial coverage on another via a learned linear projection head (20M parameters, trained on 30 aligned traces).
Mean cross-transfer coverage is 64.5\%, with same-family transfer (r1-32b $\to$ QwQ-32B, both Qwen2.5) reaching 72.1\%.
This demonstrates that safety commitment signals share geometric structure across model families.

\begin{table}[t]
\centering
\small
\begin{tabular}{lccc}
\toprule
\textbf{Train $\to$} & \textbf{r1-8b} & \textbf{r1-32b} & \textbf{QwQ} \\
\midrule
r1-8b (Llama) & .767\textsuperscript{*} & .628 & .581 \\
r1-32b (Qwen) & .651 & .791\textsuperscript{*} & .721 \\
QwQ (Qwen) & .605 & .698 & .744\textsuperscript{*} \\
\bottomrule
\end{tabular}
\caption{Cross-architecture transfer matrix (temporal $W$=15 coverage). \textsuperscript{*}In-distribution. Same-family transfer achieves $>$70\%; cross-family transfer achieves $>$58\%.}
\label{tab:transfer}
\end{table}

% ========================================================================
% 6. ABLATIONS
% ========================================================================
\section{Ablation Studies}

\subsection{Window Size}

Table~\ref{tab:window} reports coverage as a function of window size $W$ across all architectures.
Coverage increases monotonically from $W$=1 (equivalent to static) to $W$=15, then saturates or slightly declines at $W$=25 due to overfitting.
The saturation point $W$=15 is identical across all models, indicating that the temporal scale of safety reasoning decline ($\sim$15 steps) is a property of the behavioral phenomenon, not the model architecture.

\begin{table}[t]
\centering
\small
\begin{tabular}{ccccc}
\toprule
$W$ & \textbf{r1-8b} & \textbf{r1-32b} & \textbf{QwQ} & \textbf{Mean} \\
\midrule
1 & .465 & .442 & .395 & .434 \\
3 & .535 & .512 & .465 & .504 \\
5 & .651 & .628 & .581 & .620 \\
10 & .651 & .698 & .651 & .667 \\
15 & .767 & .791 & .744 & .767 \\
20 & .767 & .791 & .744 & .767 \\
25 & .744 & .767 & .721 & .744 \\
\bottomrule
\end{tabular}
\caption{Coverage vs.\ window size.  Saturation at $W$=15 is universal across architectures.  Decline at $W$=25 reflects overfitting from increased feature dimensionality.}
\label{tab:window}
\end{table}

\subsection{Layer Selection}
\label{sec:layer_ablation}

Table~\ref{tab:layers} reports coverage across relative layer depths.
The optimal layer falls at 35--44\% relative depth across all models: layer 14/32 for 8B models and layers 22--24/64 for 32B models.
This ``middle-layer sweet spot'' between low-level syntactic features and task-specific output representations is consistent with prior probing work and enables a simple layer selection heuristic (pick the middle third) without architecture-specific tuning.

\begin{table}[t]
\centering
\small
\begin{tabular}{lccc}
\toprule
\textbf{Depth} & \textbf{r1-8b} & \textbf{r1-32b} & \textbf{QwQ} \\
\midrule
10\% & .442 (L3) & .395 (L6) & .372 (L6) \\
25\% & .605 (L8) & .581 (L16) & .558 (L16) \\
35\% & .721 (L11) & \textbf{.814} (L22) & .744 (L22) \\
44\% & \textbf{.767} (L14) & .791 (L28) & \textbf{.791} (L28) \\
50\% & .744 (L16) & .767 (L32) & .767 (L32) \\
75\% & .628 (L24) & .651 (L48) & .628 (L48) \\
\bottomrule
\end{tabular}
\caption{Temporal $W$=15 coverage by relative layer depth.  Bold = best per model.  The optimal region (35--44\% depth) generalizes across architectures.}
\label{tab:layers}
\end{table}

\subsection{GRU Upper Bound}

Table~\ref{tab:gru} compares the GRU sequential probe against temporal LR on 32B models.
The GRU (1.62M parameters) does not outperform temporal LR (76.8K parameters) on any model---in fact, it underperforms by 2.3pp due to overfitting at current data scales (35\% overfitting rate on held-out validation).
This confirms that method simplicity is the correct design choice given the data regime: 209 traces do not support learning nonlinear temporal patterns beyond what linear probes capture.

\begin{table}[t]
\centering
\small
\begin{tabular}{llcccc}
\toprule
\textbf{Model} & \textbf{Method} & \textbf{Cov} & \textbf{Acc} & \textbf{Lead} & \textbf{Params} \\
\midrule
\multirow{2}{*}{r1-32b} & Temp LR & .791 & .851 & $-$1.32 & 76.8K \\
& GRU-256 & .767 & .842 & $-$0.6 & 1.62M \\
\midrule
\multirow{2}{*}{QwQ-32B} & Temp LR & .744 & .839 & $-$0.87 & 76.8K \\
& GRU-256 & .721 & .831 & $-$0.4 & 1.62M \\
\bottomrule
\end{tabular}
\caption{GRU vs.\ temporal LR on 32B models.  GRU's 21$\times$ parameter increase yields no coverage improvement, confirming that the bottleneck is data scale, not model capacity.}
\label{tab:gru}
\end{table}

\subsection{16-Label Fine-Grained Supervision}

Table~\ref{tab:16label} evaluates whether replacing binary labels with the full 16-label HarmThoughts taxonomy improves detection.
Fine-grained supervision provides marginal improvement (+2.4pp coverage, +0.3 steps lead time) but does not change the fundamental findings.
The modest gain confirms that the bottleneck is the temporal information capacity of hidden states, not label granularity.

\begin{table}[t]
\centering
\small
\begin{tabular}{lcccc}
\toprule
\textbf{Model} & \textbf{Bin Cov} & \textbf{16L Cov} & \textbf{Bin Lead} & \textbf{16L Lead} \\
\midrule
r1-8b & .767 & .791 & $-$1.09 & $-$0.7 \\
r1-32b & .791 & .814 & $-$1.32 & $-$0.9 \\
\bottomrule
\end{tabular}
\caption{Binary vs.\ 16-label supervision (temporal $W$=15).  Fine-grained labels offer marginal improvement, confirming that hidden-state temporal capacity---not supervision quality---is the bottleneck.}
\label{tab:16label}
\end{table}

\subsection{Statistical Validation}

Table~\ref{tab:bootstrap} reports 95\% bootstrap confidence intervals (1000 resamples) for key metrics.
All temporal-vs-static coverage deltas exclude zero (minimum lower bound +18.6pp), confirming statistical significance despite the moderate sample size.
Lead time CIs are centered negative with upper bounds near zero.

\begin{table}[t]
\centering
\small
\begin{tabular}{lccc}
\toprule
\textbf{Metric} & \textbf{r1-8b} & \textbf{r1-32b} & \textbf{QwQ} \\
\midrule
Static Cov & [.35, .58] & [.33, .56] & [.28, .51] \\
Temp Cov & [.67, .86] & [.70, .88] & [.65, .84] \\
$\Delta$ & [.21, .49] & [.23, .51] & [.23, .51] \\
Lead Time & [$-$2.8, 0.6] & [$-$3.1, 0.5] & [$-$2.3, 0.7] \\
\bottomrule
\end{tabular}
\caption{Bootstrap 95\% CIs (1000 resamples). All $\Delta$ CIs exclude zero.  Wide intervals reflect sample size (209 traces) but effect sizes are large relative to uncertainty.}
\label{tab:bootstrap}
\end{table}

% ========================================================================
% 7. DISCUSSION
% ========================================================================
\section{Discussion}

\paragraph{Why does method simplicity win?}
Our finding that temporal window LR matches or exceeds GRU and attention-based probes may seem counterintuitive.
The explanation is data-regime-dependent: with 209 training traces ($\sim$10K steps), a linear model with $W \times d$ features (76.8K parameters) is at the edge of the data-sufficient regime, while a GRU with 1.62M parameters overfits.
We expect complex temporal models to eventually outperform LR as datasets scale, but the current finding has practical value: it means safety monitoring can be deployed with minimal infrastructure.

\paragraph{Practical implications of negative lead time.}
While pre-textual detection remains out of reach, the temporal probe's near-zero lead time (mean $-$0.87 to $-$1.32 steps) is practically useful.
Detection at or shortly after the commitment point enables intervention strategies: (1)~post-hoc audit of reasoning traces before output delivery, (2)~real-time monitoring that halts generation within 1--2 steps of commitment, and (3)~triggering safety re-reasoning when commitment is detected.

\paragraph{Universality of safety signal geometry.}
The consistent failure of Mahalanobis distance across architectures, combined with successful cross-architecture transfer, suggests that safety commitment signals occupy a specific geometric niche---minor principal components that carry low variance but high discriminative value.
This geometry appears to be a universal property of how reasoning models encode safety state, independent of architecture family or scale.

% ========================================================================
% 8. LIMITATIONS
% ========================================================================
\section{Limitations}

\paragraph{Sample size.} Our evaluation uses 209 traces from HarmThoughts.  While bootstrap CIs confirm statistical significance of core findings, wider confidence intervals ($\pm$9--12pp) limit precision of coverage estimates.  Replication on larger datasets is needed.

\paragraph{Model coverage.} We evaluate four models from two architecture families.  Proprietary reasoning models (OpenAI o1/o3, Claude) with different training methodologies may exhibit different safety signal geometry.

\paragraph{Commitment point definition.} Our binary commitment label assumes a single, permanent transition.  Some traces may exhibit multiple commitment-like transitions or gradual commitment that is not well-captured by a point definition.

\paragraph{Probe access assumption.} Hidden-state monitoring requires white-box access to model internals, limiting deployment to first-party safety teams or open-weight models.

% ========================================================================
% 9. CONCLUSION
% ========================================================================
\section{Conclusion}

We present the first cross-architecture empirical study of hidden-state safety monitoring in reasoning models.
Across four models, two architecture families, and two scales, we establish that:
(1)~static probes fail despite strong per-dimension safety signal, because the signal concentrates in minor principal components;
(2)~a simple temporal window probe achieves oracle-ceiling coverage with a universal window size ($W$=15), requiring no architecture-specific tuning;
(3)~pre-textual detection is fundamentally limited by the causal structure of hidden-state representations, which encode current safety state but not future trajectory.
Cross-architecture transfer demonstrates shared geometric structure of safety signals across model families.
These findings provide both a practical monitoring tool (temporal window probes) and principled boundaries (lead time limitations) for the emerging challenge of reasoning model safety.

\bibliography{custom}

\end{document}
